% =====================================================================
%  APÊNDICES
%  (Material próprio, demasiado extenso para o corpo do relatório.)
% =====================================================================
\chapter{Principais \emph{Endpoints} da API}
\label{ap:endpoints}

A \Cref{tab:endpoints} resume os principais \emph{endpoints} REST expostos
pelo \emph{backend} (prefixo \texttt{/api/v1/}).

\begin{longtable}{@{}p{1.4cm}p{6.6cm}p{5.4cm}@{}}
\caption{Principais \emph{endpoints} da API.}\label{tab:endpoints}\\
\toprule
\textbf{Método} & \textbf{Rota} & \textbf{Descrição} \\
\midrule
\endfirsthead
\toprule
\textbf{Método} & \textbf{Rota} & \textbf{Descrição} \\
\midrule
\endhead
GET    & \texttt{/healthz}                       & Sonda de saúde profunda. \\
POST   & \texttt{/upload} · \texttt{/upload/batch} & Carregamento de fotografias. \\
GET    & \texttt{/media}                         & Listar \emph{media}. \\
GET    & \texttt{/media/\{id\}/file}             & Servir ficheiro (auth por \emph{header} ou \texttt{?token=}). \\
POST   & \texttt{/genealogy/upload}              & Importar GEDCOM. \\
GET    & \texttt{/genealogy/persons}             & Listar pessoas. \\
GET    & \texttt{/timeline}                      & Linha temporal. \\
GET    & \texttt{/narrative/templates}           & Listar \emph{templates}. \\
POST   & \texttt{/narrative/generate}            & Gerar narrativa (\texttt{sync}/\texttt{background}). \\
GET    & \texttt{/narrative/stories}             & Listar histórias (favoritas primeiro). \\
PATCH  & \texttt{/narrative/stories/\{id\}}       & Editar história (título/texto/favorita). \\
POST   & \texttt{/narrative/stories/\{id\}/regenerate} & Regenerar a narrativa com \emph{feedback}. \\
POST   & \texttt{/multimedia/generate/\{id\}}    & Gerar vídeo a partir de uma história. \\
GET    & \texttt{/multimedia/videos}             & Listar vídeos. \\
GET    & \texttt{/multimedia/subtitle/\{ficheiro\}} & Servir a faixa de legendas \texttt{.vtt}. \\
GET/POST & \texttt{/projects} · \texttt{/projects/\{id\}} & CRUD de projetos. \\
GET    & \texttt{/tasks} · \texttt{/tasks/\{id\}} & Estado de tarefas assíncronas. \\
\bottomrule
\end{longtable}

\chapter{Configuração de \emph{Deployment}}
\label{ap:deploy}

Resumo da topologia de \emph{deployment} (ver \Cref{sec:impl_deploy}):

\begin{itemize}[leftmargin=*]
  \item \textbf{Vercel} — \emph{frontend} (SPA), \emph{build} Vite,
  reescrita SPA.
  \item \textbf{Render} — \emph{backend} via Docker (FFmpeg, libmagic,
  Tesseract), \texttt{render.yaml}, sonda \texttt{/healthz},
  \texttt{CELERY\_ENABLED=False} e ChromaDB desativado.
  \item \textbf{Supabase} — Auth (JWKS), PostgreSQL (asyncpg), Storage.
\end{itemize}

\noindent\emph{Nota:} as variáveis sensíveis (chaves Supabase, Groq e
Gemini) são definidas no painel do Render e nunca versionadas no
repositório.
